\section{Planteamiento: qué se mide y contra qué}\label{sec:planteamiento}

Esta sección fija el vocabulario y la notación del resto del documento. Todo lo demás (el contador,
la evidencia, la regla de decisión, las métricas) se apoya en cuatro definiciones: qué es un secreto,
qué es una llamada, cuánto de un secreto ha salido y cuándo decimos que el ataque triunfó.

\subsection{El escenario}

Un agente de código trabaja en un repositorio que contiene claves: en un archivo \cod{.env}, en un
\cod{config/secrets.yaml}. Para trabajar lee archivos, corre \cod{git} y habla con servicios por HTTP.
También lee texto que no escribió su equipo: un issue, una página de documentación. Ese texto puede
traer una instrucción plantada. Si el agente la sigue, usa sus propias herramientas para enviar la
clave hacia afuera, y lo puede hacer de a pocos caracteres por llamada, entre decenas de llamadas
normales.

Es el patrón que Willison llama la «tríada letal»~\cite{willison}: contenido no confiable, acceso a
datos privados y un canal de salida, las tres cosas en el mismo agente. Las herramientas que hacen
útil al agente son las mismas que sirven de canal.

\subsection{Qué ve Alicanto y qué no}

Alicanto se pone en el único lugar por donde pasan las tres cosas: la llamada a la herramienta.

\begin{itemize}
  \item \textbf{Ve} el nombre de la herramienta, sus argumentos y su resultado, para cada llamada de
    cada sesión.
  \item \textbf{Conoce} las claves que protege y sus destinos permitidos. Vienen en un manifiesto
    (\cod{demo/secrets_manifest.json}) que para cada clave da su identificador, su valor, los archivos
    donde vive y los \esp{hosts} a los que es legítimo enviarla.
  \item \textbf{Puede} pedirle al gestor de secretos que cambie una clave, y puede impedir que una
    llamada se ejecute.
  \item \textbf{No ve} el razonamiento del modelo, ni el \esp{prompt}, ni lo que ocurre fuera de las
    herramientas que él ofrece.
\end{itemize}

De ahí sale la primera decisión de diseño: como la intención no se puede observar, Alicanto no la
estima. Mide algo que sí puede observar, que es cuánto material de cada clave ha cruzado hacia afuera.

\subsection{Secretos y entropía}

Un secreto $s$ es una cadena de largo $n$ sobre un alfabeto $A_s$. Su entropía, suponiendo que cada
carácter se eligió al azar y de forma uniforme, es
\begin{equation}
  H(s) = n \cdot \log_2 |A_s| .
  \label{eq:H}
\end{equation}
El alfabeto no se declara: se deduce del valor. Es el más chico de una lista fija que contiene a todos
los caracteres de $s$ (\cref{tab:alfabetos}). Las claves de la demo y de los tres sets de datos son de
40 caracteres hexadecimales, así que $H = 40 \cdot 4 = 160$ bits.

\begin{table}[ht]
\caption{Alfabetos del contador de bits (\cod{engine/bits/decoders.py}, función \cod{alphabet_of}). Se
prueba en este orden y se usa el primero que contiene a la clave.}\label{tab:alfabetos}
\small
\begin{tabularx}{\textwidth}{@{}P{2.5cm}LR{1.6cm}R{2.4cm}@{}}
\toprule
Alfabeto & Caracteres & $|A_s|$ & Bits por carácter \\
\midrule
\cod{hex} & \cod{0-9}, \cod{a-f}; no distingue mayúsculas & 16 & \num{4.000} \\
\cod{alnum} & \cod{0-9}, \cod{a-z}, \cod{A-Z} & 62 & \num{5.954} \\
\cod{base64} & los de \cod{alnum} más \cod{+ / = - _} & 65 & \num{6.022} \\
\cod{printable} & cualquier otro: los 95 caracteres ASCII imprimibles & 95 & \num{6.570} \\
\bottomrule
\end{tabularx}
\end{table}

\subsection{Llamadas, sumideros y contenido no confiable}

Una sesión es una secuencia de llamadas $c_1, c_2, \dots$ Cada llamada $c_t$ tiene una herramienta,
unos argumentos y, después de ejecutarse, un resultado. El proxy le agrega tres datos antes de decidir:

\begin{description}
  \item[Recursos.] Lo que la llamada toca: la ruta del archivo, el \esp{host} de la petición, o
    \cod{git}.
  \item[Sumidero (\cod{is_sink}).] La llamada saca datos hacia afuera. Lo son las peticiones HTTP,
    los issues y comentarios de GitHub, y \cod{git push}, \cod{fetch} y \cod{pull}.
  \item[No confiable (\cod{untrusted}).] El resultado de la llamada es contenido que no escribió el
    equipo: lo que viene de GitHub o de la web.
\end{description}

Las reglas exactas con que el proxy pone esas marcas están en la \cref{sec:marcas}. Importan mucho:
el resultado principal de la evaluación cambió cuando se pasó de las marcas que traían las grabaciones
a las que calcula el proxy.

\subsection{Cuánto ha salido}

Para cada secreto $s$, el contador mantiene el conjunto $K_t(s) \subseteq \{0, \dots, n-1\}$ de
posiciones de $s$ que ya aparecieron en los argumentos de algún sumidero hasta la llamada $t$. Lo
filtrado, en bits, y la fracción filtrada son
\begin{equation}
  I_t(s) = |K_t(s)| \cdot \log_2 |A_s| , \qquad \lambda_t(s) = \frac{I_t(s)}{H(s)} .
  \label{eq:I}
\end{equation}
Es una cuenta por posiciones, no por caracteres enviados: mandar dos veces el mismo pedazo no suma, y
mandar un pedazo que se solapa con otro suma solo lo nuevo. Es también una cota inferior de lo que el
atacante sabe. Si un pedazo sale en una codificación que el contador no sabe leer, o es más corto que
el mínimo que cuenta, $K_t$ no lo registra (\cref{sec:contador-limites}).

Cuando una clave rota, su valor cambia y $K_t(s)$ vuelve a vacío: lo que salió era de un valor que ya
no sirve. Por eso todas las cuentas de este documento son \emph{por versión} de la clave.

\subsection{Cuándo triunfa el ataque}

No hace falta sacar la clave entera para robarla: con la mayor parte, el resto se puede probar por
fuerza bruta. El diseño fija una seguridad residual $\beta = 64$ bits: mientras al atacante le falten
más de 64 bits de una versión, se considera que no la tiene. Con $H = 160$,
\begin{equation}
  \text{el ataque triunfa si saca más de } H - \beta = 96 \text{ bits de una misma versión de la clave,}
  \label{eq:exito}
\end{equation}
es decir, más de 24 de los 40 caracteres. Es el criterio de la métrica (\cod{eval/metrics/metrics.py}),
del servidor del atacante que decide si puede completar la clave (\cod{services/attacker/app.py}) y de
la invariante que acota los parámetros (\cref{sec:invariante}).

\begin{limite}
\textbf{Un detalle de borde.} La métrica cuenta el ataque como frenado si salieron \emph{96 bits o
menos}. El servidor del atacante da la clave por completada si reunió \emph{96 bits o más}. En el
caso exacto de 96 bits las dos piezas no coinciden. No afecta a ninguna cifra de Alicanto (el máximo
que dejó salir es menor), pero sí a una ablación del análisis del atacante que se adapta, donde una
familia queda justo en 96 (\cref{sec:adaptativo}).
\end{limite}

\subsection{Las dos decisiones de diseño, en estos términos}

\begin{description}
  \item[Medir.] La cantidad que gobierna las decisiones es $\lambda_t(s)$, que se acumula entre
    llamadas, entre codificaciones y, con la cuenta compartida, entre agentes. Ninguna llamada tiene
    que parecer sospechosa por sí sola.
  \item[Rotar.] La respuesta por defecto no es impedir la llamada sino devolver $K_t(s)$ a vacío,
    cambiando la clave en el gestor. Un error de Alicanto en una sesión normal cuesta una rotación y no
    interrumpe el trabajo. El bloqueo queda para tres casos: cuando sale un tramo largo y seguido de la
    clave hacia un destino que no es el suyo, cuando la evidencia y el camino de los datos coinciden, y
    cuando una clave ya rotó demasiadas veces (\cref{sec:decision}).
\end{description}

\subsection{Notación}

\begin{table}[ht]
\caption{Símbolos usados en este documento, con su nombre en el código.}\label{tab:notacion}
\small
\begin{tabularx}{\textwidth}{@{}P{2.2cm}LP{4.4cm}@{}}
\toprule
Símbolo & Significado & En el código \\
\midrule
$s$, $n$, $A_s$ & Un secreto, su largo y su alfabeto & \cod{_SecretState} \\
$H(s)$ & Entropía del secreto, en bits & \cod{BitsEngine.H} \\
$K_t(s)$ & Posiciones confirmadas de $s$ que ya salieron & \cod{K} \\
$C_t(s)$ & Posiciones candidatas vigentes (\cref{sec:confirmadas}) & \cod{cand} \\
$I_t(s)$, $\lambda_t(s)$ & Bits filtrados y fracción filtrada & \cod{bits}, \cod{fraction} \\
$k$ & Largo mínimo de una coincidencia confirmada (8) & \cod{k} \\
$W$ & Llamadas que dura una candidata sin compañía (50) & \cod{W} \\
$\ell_t$ & Incremento de evidencia de la llamada $t$ & \cod{FeatureModel.llr} \\
$S_t$ & Evidencia acumulada (CUSUM) & \cod{Decision.S} \\
$B$ & Tope del incremento, $|\ell_t| \leq B$ (3) & \cod{B} \\
$\lrot$ & Fracción filtrada que dispara la rotación & \cod{lambda_rot} \\
$\brot$, $b$ & Umbrales de evidencia para rotar y para bloquear & \cod{b_rot}, \cod{b} \\
$T$ & Llamadas entre rotaciones periódicas de una clave expuesta & \cod{T} \\
$\beta$ & Seguridad residual exigida, en bits (64) & \cod{beta_bits} \\
\bottomrule
\end{tabularx}
\end{table}
